% Book-wide notation and conventions. Each chapter's own notation box adds the letters it uses locally.
The symbols below keep the same meaning in every chapter. Letters that a chapter uses for something
else are listed in the last part of the table and are defined again in that chapter's notation box.

\begingroup\small
\setlength{\LTpre}{6pt}\setlength{\LTpost}{0pt}
\begin{longtable}{@{}>{\raggedright\arraybackslash}p{0.24\textwidth}>{\raggedright\arraybackslash}p{0.70\textwidth}@{}}
\toprule
Symbol & Meaning\\
\midrule
\endhead
\bottomrule
\endfoot
\multicolumn{2}{@{}l}{\textbf{Probability}}\\
$\mathbb P$, $Q$ & Real-world probability and a risk-neutral measure; $Q^T$ is the $T$-forward measure.\\
$\E$, $\E_Q$, $\E_{\mathbb P}$ & Expectation, and expectation under the measure in the subscript.\\
$\E_t$ & Expectation given the information available at $t$.\\
$\Var$, $\Cov$, $\Corr$ & Variance, covariance and correlation.\\
$\Phi$, $\varphi$ & Standard normal distribution and density functions.\\
$\N(\mu,\sigma^2)$ & Normal law with mean $\mu$ and variance $\sigma^2$.\\
$\mathbf 1\{A\}$, $x^+$ & Indicator of the event $A$, and $\max(x,0)$.\\
$W_t$ & Standard Brownian motion ($Z_t$ in Chapter~11).\\[3pt]
\multicolumn{2}{@{}l}{\textbf{Time}}\\
$t$, $T$ & Calendar time in years, and a maturity or horizon.\\
$\tau$ & Time to expiry or term, $T-t$, in Chapters~4, 5, 9 and~12; a stopping time in Chapter~3 and
\cref{sec4:american}, and a default time in Chapter~6.\\[3pt]
\multicolumn{2}{@{}l}{\textbf{Returns and rates}}\\
$R_t$ & Simple return over a period, $(V_t+D_t-V_{t-1})/V_{t-1}$.\\
$r_t$ & Log return $\ln(1+R_t)$ in Chapter~1; a daily return in Chapter~8, a monthly excess return
in Chapter~9 and a true, unsmoothed return in Chapter~10.\\
$r$ & A continuously compounded riskless rate in the pricing chapters, 3 to~6; $r_f$ is the riskless
rate in Chapter~7.\\
$y$ & A yield. Its compounding is stated with it; see the conventions below.\\
$\delta$ & Continuous dividend or income yield of an asset (Chapters~4 and~5).\\
$\sigma$ & Volatility, annual unless the period is given.\\
$\mu$ & Mean return or drift.\\
$\theta$ & Sharpe ratio in Chapters~7, 8, 9 and~12; the market price of risk in Chapters~3 and~5.\\
$\gamma_3$, $\gamma_4$ & Skewness and \emph{excess} kurtosis. Every kurtosis printed in the book is excess
kurtosis, which is zero for a normal law; Chapter~1 writes $S$ and $K_e$ for the same quantities.\\[3pt]
\multicolumn{2}{@{}l}{\textbf{Funds, fees and smoothing}}\\
$h$ & Hurdle rate. A \emph{soft} hurdle pays the manager a share of the whole gain once the hurdle is
passed; a \emph{hard} hurdle pays a share of the gain above it only.\\
$c$, $k$ & Carried interest rate and catch-up rate (Chapters~2 and~11).\\
$m$, $p$, $f$ & Management fee rate, and incentive fee rate ($p$ in Chapter~2, $f$ in Chapter~9).\\
$M_t$ & High-water mark: the highest value per unit after fees at an earlier crystallisation date, or
the starting value if that is higher.\\
$\phi$ & Smoothing parameter of reported returns, $r^*_t=(1-\phi)r_t+\phi r^*_{t-1}$ (Chapters~1, 7, 8
and~10; Chapter~1 writes $R$ for $r$).\\
$C_t$, $D_t$ & Contributions to and distributions from a fund (Chapter~11; Chapter~2 writes $C_i$ and $L_j$).\\
$\mathit{NOI}$, $V$, $c$ & Net operating income, value and cap rate of a property (Chapter~10; $q$ in the
case study of Chapter~6).\\[3pt]
\multicolumn{2}{@{}l}{\textbf{Letters with local meanings}}\\
$\lambda$ & Hazard rate (Chapter~6), jump intensity $\lambda_J$ (Chapter~5), storage rate (Chapter~4),
EWMA decay or drawdown rate (Chapter~8), loan to value (Chapter~10), debt share (Chapter~11),
penalty (Chapter~12).\\
$\tau$ & Also the date of the last contribution (Chapter~2), the prior scale of Black-Litterman
(Chapter~7), the Omega threshold (Chapter~8) and a tax rate (Chapter~11).\\
$\alpha$ & Regression intercept (Chapter~7) and a confidence level (Chapter~8).\\
$\beta$ & Regression slope or market beta; the rate at which cap rates follow yields in the case study of
Chapter~6.\\
$q$ & A risk-neutral probability (Chapter~4); a number of lags or a vector of views (Chapter~7);
a cap rate (Chapter~6).\\
$\eta$, $y$ & Convenience yield: $\eta$ in Chapter~4, $y$ in Chapters~9 and~10.\\
\end{longtable}
\endgroup

\paragraph{Compounding and units}
Rates are quoted per year. Each chapter keeps one convention for each kind of rate and states the
conversion where two meet:
\begin{itemize}[nosep]
  \item Chapters 1 and 2: effective annual rates unless another compounding is named; an IRR is an
        effective annual rate on dated flows, and $\ln(1+R)$ converts to continuous compounding.
  \item Chapters 3 to 5: continuously compounded rates; Chapter~4 also uses simple money market rates
        and yields compounded $m$ times a year, $D(t)=(1+y/m)^{-mt}$.
  \item Chapter 6: bond yields compounded twice a year, Treasury bill yields simple on an actual/365
        basis, continuously compounded rates inside the short-rate models, loan rates compounded
        monthly in the private credit screen, and half-yearly in the mortgage case study, as the
        Interest Act requires a Canadian mortgage rate to be stated.
  \item Chapters 7 to 9 and 12: monthly or daily returns, simple unless they are called log returns,
        as the forecasting target of Chapter~12 is; annualised figures say how they were annualised.
  \item Chapter 10: log returns in the smoothing model and regressions, simple annual rates in the
        valuation and loan formulas, mortgage payments compounded monthly, continuous rates in the
        forest and storage sections.
  \item Chapter 11: annual rates compounded once a year unless marked as continuous.
\end{itemize}
Percentages in the text are of the quantity named; a basis point is 0.01 percentage point. Money
amounts carry their currency where it matters and are otherwise per 100 of a stated base.
